% 针孔投影的几何关系；示意深度候选，不按物理长度绘制。
\begin{tikzpicture}[font=\figurefont\small,x=1cm,y=1cm,
 ray/.style={draw=BookBlue,line width=0.9pt},
 alt/.style={draw=BookAmber,dashed,line width=0.9pt}]
 \coordinate (a) at (0,0);
 \coordinate (b) at (7,0);
 \coordinate (p) at (2,3);
 \coordinate (q) at (3,4.5);
 \draw[ray] (a)--(q);
 \draw[ray] (b)--(p);
 \draw[alt] (b)--(q);
 \draw[draw=BookInk,line width=1.2pt] (-0.3,1)--(2.2,1);
 \draw[draw=BookInk,line width=1.2pt] (4.5,1)--(7.5,1);
 \fill[BookInk] (a) circle(1.1mm);
 \fill[BookInk] (b) circle(1.1mm);
 \fill[BookTeal] (p) circle(1.1mm);
 \draw[BookAmber,fill=white,line width=0.9pt] (q) circle(1.1mm);
 \fill[BookBlue] (0.667,1) circle(1mm);
 \fill[BookBlue] (5.333,1) circle(1mm);
 \draw[BookAmber,fill=white,line width=0.9pt] (6.111,1) circle(1mm);
 \node[below] at (a) {参考相机};
 \node[below] at (b) {源相机};
 \node[above right] at (0.667,1) {同一参考像素};
 \node[below] at (5.333,0.9) {$u'_1$};
 \node[below] at (6.111,0.9) {$u'_2$};
 \node[left] at (p) {深度$d_1$};
 \node[right] at (q) {深度$d_2$};
 \node[right] at (7.5,1) {源图平面};
 \draw[draw=BookMuted,<->,line width=0.7pt] (0,-0.8)--(7,-0.8)
   node[midway,below] {已知相机相对位置};
\end{tikzpicture}
